\chapter*{Introduction générale}
\addcontentsline{toc}{chapter}{Introduction générale}
\markboth{Introduction générale}{}

Un robot mobile autonome n'est exploitable que si un opérateur peut lui confier une tâche
et suivre son exécution. Dans un entrepôt ou un hôpital, cet opérateur ne travaille pas
dans un terminal ROS : il utilise une interface de supervision qui affiche la position des
robots, lui permet de désigner une destination et l'informe du résultat. La qualité de
cette interface dépend moins de son apparence que de la fidélité des informations
affichées : une position fausse, un ordre envoyé deux fois ou un succès annoncé à tort
peuvent conduire l'opérateur à une mauvaise décision.

Le présent stage s'inscrit dans le projet AMR-X, une plateforme de robot mobile modulaire
développée sous ROS~2. Au début du stage, le projet disposait déjà d'une simulation Gazebo
du robot dans un entrepôt, d'une pile de navigation Nav2 configurée et d'un dashboard web
de supervision. Ces trois briques fonctionnaient séparément : le dashboard affichait une
carte du monde mais ne pouvait ni connaître la position réelle du robot simulé, ni lui
envoyer une destination. L'objectif du stage est de relier ces briques afin que
l'opérateur puisse, depuis le dashboard, choisir une destination, commander le robot
simulé dans Gazebo et voir sa position et le résultat de la navigation dans le navigateur.

La démarche adoptée consiste d'abord à analyser l'existant et à définir un protocole de
communication sûr entre le navigateur et Nav2, puis à développer une passerelle ROS~2 et
les composants correspondants du dashboard, et enfin à valider l'ensemble dans la
simulation par des essais mesurés. Une attention particulière est portée à la cohérence
entre le repère de la carte Nav2 et celui du monde Gazebo, qui s'est révélée déterminante
lors des essais.

Ce rapport comporte quatre chapitres. Le premier présente le contexte du projet AMR-X,
l'existant, la problématique et le cahier des charges. Le deuxième décrit la conception de
la solution : architecture, protocole de communication, machine d'états et
transformation des repères. Le troisième détaille la réalisation de la passerelle ROS~2, de
l'interface web et de l'outillage de lancement, ainsi que les problèmes rencontrés. Le
quatrième présente le protocole de validation et les résultats mesurés en simulation.
